\section*{Chapitre \ref{ch:g12:past-climates} --- Lire les climats du passé}

\begin{solution}{exo:g12:past-climates:1}
(a) un \omterm{def:g12:past-climates:proxy}{indicateur climatique}~; (b) une \omterm{def:g12:past-climates:ice-core}{carotte de glace}~; (c) un
\omterm{def:g12:past-climates:glacial-cycle}{interglaciaire}~; (d) l'\omterm{def:g12:past-climates:pollen}{analyse pollinique}~; (e) le \omterm{def:g12:past-climates:lgm}{dernier maximum
glaciaire}.
\end{solution}

\begin{solution}{exo:g12:past-climates:2}
(a) La glace polaire~: le dioxyde de carbone et le méthane de l'air de
l'époque. (b) Les sédiments lacustres ou la tourbe~: la végétation, ici des
arbres qui ont besoin d'étés chauds. (c) Les sédiments océaniques~: le
volume de glace sur les continents et la température de l'\omterm{def:g10:layered-ocean:thermocline}{océan profond}.
(d) Les arbres~: les conditions de croissance, année par année.
\end{solution}

\begin{solution}{exo:g12:past-climates:3}
(a) $-50 - (-120) = \qty{70}{m}$ en \num{7000} ans~: $\qty{70000}{mm}/\num{7000}
= \qty{10}{mm}$ par an. (b) $10/3.7 \approx 2.7$~: la montée moyenne de la
déglaciation était environ 2.7 fois plus rapide que celle de 2006--2018.
\end{solution}

\begin{solution}{exo:g12:past-climates:4}
Un siècle compte dix décennies~: $\qty{2.0}{\degreeCelsius}$ par siècle.
Mille ans comptent cent décennies~: $\qty{20}{\degreeCelsius}$ par millier
d'années.
\end{solution}

\begin{solution}{exo:g12:past-climates:5}
(a) \num{23000} ans~: la précession de l'axe~; \num{42000} ans~:
l'inclinaison de l'axe~; environ \num{100000} ans~: la forme de l'orbite.
(b) $\num{1000000}/\num{41000} \approx 24$ cycles~; $\num{1000000}/\num{100000}
= 10$ cycles.
\end{solution}

\begin{solution}{exo:g12:past-climates:6}
(a) $266.6 + 7.83 \times (-9.4) \approx 266.6 - 73.6 = \qty{193}{ppm}$~:
seulement \qty{3}{ppm} au-dessus des \qty{190}{ppm} mesurés. (b) $266.6 + 7.83 \times 5.5
\approx \qty{310}{ppm}$, \qty{11}{ppm} au-dessus de la plus haute valeur
mesurée~: le point le plus chaud se trouve sous la droite. (c) $\Delta T = (427 - 266.6)/7.83 \approx
\qty{20}{\degreeCelsius}$. La droite ne décrit que la façon dont les deux
grandeurs ont varié ensemble, au cours de cycles naturels, dans le domaine
des données (\qty{-11}{\degreeCelsius} à \qty{+6}{\degreeCelsius}, 174 à
\qty{299}{ppm})~; ce n'est pas une loi, et l'on ne peut pas y lire une
valeur située très loin de ce domaine.
\end{solution}

\begin{solution}{exo:g12:past-climates:7}
(a) Les herbacées, \qty{84}{\percent}~: un paysage ouvert presque sans
arbres, froid. (b) Le chêne (\qty{42}{\percent}) et le noisetier
(\qty{34}{\percent})~: une forêt de feuillus, avec des étés chauds. (c) Le
bouleau s'étend le premier, puis le pin~; le noisetier et le chêne arrivent
les derniers. (d) Les arbres ne s'étendent qu'à la vitesse à laquelle leurs
graines voyagent et leurs jeunes plants grandissent, à partir des lieux où
ils ont survécu au froid~; un sol doit aussi se former. La forêt peut être
en retard sur le climat.
\end{solution}

\begin{solution}{exo:g12:past-climates:8}
(a) \qty{300}{cm}, soit \qty{3000}{mm}, en \num{8000} ans~: $3000/8000 =
\qty{0.375}{mm}$ par an. (b) $\num{6000} + (470 - 300)/300 \times \num{8000}
\approx \num{6000} + \num{4530} \approx \num{10500}$ ans. (c) $10/0.375
\approx 27$ ans.
\end{solution}

\begin{solution}{exo:g12:past-climates:9}
(a) $7.16 - 4.11 = 3.05$ millions de kilomètres carrés, et $3.05/4.11 \approx
\qty{74}{\percent}$ de terres en plus. (b) Oui~: au \omterm{def:g12:past-climates:lgm}{dernier maximum
glaciaire}, la mer se trouvait 125 à \qty{134}{m} sous son niveau actuel,
bien plus que les \qty{36}{m} qui relient Bornéo à l'Asie sur la carte.
(c) La carte utilise les fonds marins actuels, qui ne sont pas restés
exactement les mêmes depuis (des sédiments s'y sont déposés, et le sol a
un peu bougé)~; la grille est aussi grossière, un point tous les
0.1~degré, et le bas niveau marin n'est connu qu'à une dizaine de mètres
près.
\end{solution}

\begin{solution}{exo:g12:past-climates:10}
(a) La valeur du \omterm{def:g12:past-climates:lgm}{dernier maximum glaciaire}, \qty{4.92}{\permil}, est la plus
élevée. C'est ce que l'on attend~: les grandes calottes glaciaires de
l'époque stockaient une eau pauvre en oxygène~18, ce qui laissait l'océan,
et les coquilles qui y croissaient, plus riches en oxygène~18~; l'eau
profonde était aussi plus froide. (b) Pour que les époques froides et
englacées soient en bas et les époques chaudes en haut, comme sur une
courbe de température.
\end{solution}

\begin{solution}{exo:g12:past-climates:11}
(a) Une corrélation montre que les deux enregistrements montent et
descendent ensemble, près d'une droite~; elle ne montre ni lequel a causé
l'autre, ni même que l'un a causé l'autre. (b) Un climat plus chaud aurait
pu libérer du dioxyde de carbone dans l'air~; ou une troisième cause, comme
les variations de l'orbite de la Terre, aurait pu faire varier les deux~;
ou chacun aurait pu pousser l'autre. (c) La chronologie des changements,
lequel est venu en premier, grâce à une datation précise de la glace et de
son air~; et la compréhension de la façon dont le dioxyde de carbone agit
sur le rayonnement de la Terre (\cref{ch:g12:climate-change}).
\end{solution}

\begin{solution}{exo:g12:past-climates:12}
(a) $\num{1300000}/\num{41000} \approx 32$ cycles entre 2.6 et 1.3
millions d'années~; $\num{700000}/\num{98500} \approx 7$ cycles sur les
\num{700000} dernières années. (b) \num{41000} ans~: l'inclinaison de
l'axe~; environ \num{100000} ans~: la forme de l'orbite. (c) Un cycle de
\num{41000} ans est échantillonné environ $\num{41000}/\num{2000} \approx 20$ fois,
assez pour le tracer~; un changement qui dure un siècle tombe entre deux
échantillons, ou à l'intérieur d'un seul, et ne laisse presque pas de
trace.
\end{solution}

\begin{solution}{exo:g12:past-climates:13}
(a) $192/10 \approx 19$ et $198/10 \approx 20$~: environ 20 fois plus vite.
(b) Contre \qty{4}{ppm} par siècle~: $192/4 = 48$ à $198/4 \approx 50$
fois~; contre \qty{42}{ppm}~: $192/42 \approx 4.6$ à $198/42 \approx 4.7$
fois. La hausse actuelle est environ 5 à 50 fois plus rapide que celle du
PETM. (c) La hausse vaut $2000 - 900 = \qty{1100}{ppm}$~: $1100/42 \approx 26$ siècles,
soit environ \num{2600} ans au rythme le plus rapide du PETM~;
$1100/195 \approx 5.6$ siècles, moins de 600 ans, au rythme actuel.
\end{solution}

\begin{solution}{exo:g12:past-climates:14}
(a) Dix fois 1 à \qty{1.5}{\degreeCelsius}~: 10 à \qty{15}{\degreeCelsius}
par millier d'années, soit 0.10 à \qty{0.15}{\degreeCelsius} par décennie.
(b) Le rythme mondial actuel, \qty{0.20}{\degreeCelsius} par décennie, vaut
1.3 à 2 fois le réchauffement le plus rapide au-dessus du Groenland. (c) Un
lieu isolé peut changer bien plus que la moyenne sur le globe, où les
changements des différentes régions se compensent en partie~: on compare
un rythme mondial à un rythme mondial.
\end{solution}

\begin{solution}{exo:g12:past-climates:15}
(a) \qty{14}{m} à \qty{12}{m} par millier d'années prennent $14/12 \approx
1.17$ millier d'années, environ \num{1170} ans~: il y a environ \num{9800}
ans. (b) Non~: la mer devait rester au moins \qty{36}{m} sous son niveau
actuel pour relier Bornéo à l'Asie~; une fois \qty{-36}{m} dépassé, le pont
terrestre fut noyé, et Bornéo redevint une île. (c) $36/12 = 3$ milliers
d'années, environ \num{3000} ans~; en réalité, la montée a ralenti après
il y a environ \num{8200} ans~: elle a donc pris plus longtemps.
\end{solution}

\begin{solution}{pb:g12:past-climates:1}
\textbf{1.} La glace elle-même, par ses isotopes lourds, enregistre la
température~; l'air de ses bulles enregistre le dioxyde de carbone.

\textbf{2.} $5.5 - (-10.6) = \qty{16.1}{\degreeCelsius}$.

\textbf{3.} $299 - 174 = \qty{125}{ppm}$.

\textbf{4.} Les points sont proches d'une droite~: les deux
enregistrements ont monté et baissé ensemble. Cela ne dit pas lequel a
causé l'autre.

\textbf{5.} Les enregistrements du passé sont tracés avec les époques les
plus anciennes d'un côté ou de l'autre~; lu dans le mauvais sens, un
réchauffement serait pris pour un refroidissement.

\textbf{6.} Une \omterm{def:g12:past-climates:glacial-cycle}{période glaciaire} est un intervalle froid au sein d'une
glaciation, pendant lequel les calottes glaciaires s'étendent bien au-delà
de leur taille actuelle~; un \omterm{def:g12:past-climates:glacial-cycle}{interglaciaire} est un intervalle chaud entre
deux \omterm{def:g12:past-climates:glacial-cycle}{périodes glaciaires}, avec des calottes glaciaires proches de celles
d'aujourd'hui ou plus petites.

\textbf{7.} Dans un \omterm{def:g12:past-climates:glacial-cycle}{interglaciaire}, qui a commencé il y a environ
\num{11000} ans.

\textbf{8.} $\num{800000}/\num{98000} \approx 8$~: environ huit cycles.

\textbf{9.} La forme de l'orbite.

\textbf{10.} Un rythme de \num{41000} ans, la période de l'inclinaison de
l'axe.

\textbf{11.} $-0.4 - (-9.4) = \qty{9.0}{\degreeCelsius}$.

\textbf{12.} $264 - 190 = \qty{74}{ppm}$.

\textbf{13.} \qty{6}{\degreeCelsius}~; la variation antarctique,
\qty{9.0}{\degreeCelsius}, est plus grande, environ 1.5 fois la variation
mondiale.

\textbf{14.} $\num{18000} - \num{11000} = \num{7000}$ ans.

\textbf{15.} $6/7 \approx \qty{0.86}{\degreeCelsius}$ par millier d'années.

\textbf{16.} Mille ans comptent cent décennies~: $0.86/100 \approx
\qty{0.0086}{\degreeCelsius}$ par décennie.

\textbf{17.} $0.20 \times 100 = \qty{20}{\degreeCelsius}$ par millier
d'années.

\textbf{18.} $20/0.86 \approx 23$ fois.

\textbf{19.} $20/1.5 \approx 13$ et $20/1 = 20$~: 13 à 20 fois.

\textbf{20.} Le réchauffement climatique actuel est \textbf{environ vingt
fois plus rapide} que le réchauffement de la sortie de la dernière
glaciation~: 23 fois son rythme moyen, 13 à 20 fois son rythme le plus
rapide.
\end{solution}
